% !PRINTED 306
\weberpage{306}
%<<CONT
\[
\begin{aligned}
S_3 &= S_4 + \begin{pmatrix} -1, & 1 \\ 1, & 1 \end{pmatrix} S_4 \\
S_2 &= S_3 + \begin{pmatrix} 0, & 1 \\ -1, & 0 \end{pmatrix} S_3 \\
S_1 &= S_2 + \begin{pmatrix} 1, & 0 \\ 1, & 1 \end{pmatrix} S_2 + \begin{pmatrix} 1, & 0 \\ -1, & 1 \end{pmatrix} S_2 \\
S &= S_1 + \begin{pmatrix} -1, & 0 \\ 0, & 1 \end{pmatrix} S_1.
\end{aligned}
\]

We accordingly have the composition and index series of $S$

\[
\begin{gathered}
S,\ S_1,\ S_2,\ S_3,\ S_4, \\
2 \quad 3 \quad 2 \quad 2
\end{gathered}
\]

\gesperrt{and we thus arrive here at the result that all equations
of degree 9 with linear congruence group
are metacyclic.}

\begin{center}
\S.\,79.

\gesperrt{Ternary linear congruence group for the modulus 2.}
\end{center}

Manifold interesting relations arise when
we consider the ternary linear congruence group $R$ for the modulus 2,
which can be conceived as a transitive permutation group of eight
elements, and which must contain the group of the metacyclic
equations of degree 8. By what was
proved above, it is above all a matter of considering the homogeneous
congruence group $S$, which consists of the elements

\weq{1}{\sigma = \begin{pmatrix} a, & b, & c \\ a_1, & b_1, & c_1 \\ a_2, & b_2, & c_2 \end{pmatrix}}

where the figures $a, b, c, \ldots$ are to be taken according to the modulus 2, hence $= 0$
or $= 1$, and which are compounded according to the rule given
in \S.\,37. Only such systems
of the numbers $a, b, c, \ldots$ come into consideration here whose determinant
is $= 1$, i.\,e. odd.

To determine the degree of the group $S$, observe that
the first row of $\sigma$ can have seven different forms:

\[
(1,0,0),\ (0,1,0),\ (0,0,1),\ (0,1,1),\ (1,0,1),\ (1,1,0),\ (1,1,1).
\]
